\documentclass{article}

\usepackage{amsmath,amssymb}
\usepackage{xurl}
\usepackage[hidelinks]{hyperref}
\setlength{\emergencystretch}{2em}

% Shared commands for notes in docs/paper-gaps/.

\DeclareUnicodeCharacter{2080}{\ifmmode{}_{0}\else\textsubscript{0}\fi}
\DeclareUnicodeCharacter{2081}{\ifmmode{}_{1}\else\textsubscript{1}\fi}
\DeclareUnicodeCharacter{2082}{\ifmmode{}_{2}\else\textsubscript{2}\fi}
\DeclareUnicodeCharacter{2099}{\ifmmode{}_{n}\else\textsubscript{n}\fi}

\newcommand{\Lean}{\textsc{Lean}}
\ExplSyntaxOn
\NewDocumentCommand{\leanid}{m}{
  \ifmmode
    \textsf{\detokenize{#1}}
  \else
    \group_begin:
    \urlstyle{sf}
    \tl_set:Nn \l_tmpa_tl {#1}
    \tl_replace_all:Nnn \l_tmpa_tl {\_} {_}
    \exp_args:NV \path \l_tmpa_tl
    \group_end:
  \fi
}
\ExplSyntaxOff
\providecommand{\tr}{\operatorname{tr}}
\newcommand{\tnleanIssueBase}{https://github.com/LionSR/TNLean/issues/}
\newcommand{\tnleanPrBase}{https://github.com/LionSR/TNLean/pull/}
\newcommand{\tnleanDocBase}{https://sirui-lu.com/TNLean/docs/}
\newcommand{\ghissue}[1]{\href{\tnleanIssueBase#1}{GitHub issue \##1}}
\newcommand{\ghpr}[1]{\href{\tnleanPrBase#1}{PR \##1}}
\newcommand{\doclink}[1]{\href{\tnleanDocBase#1.html}{\path{#1}}}

% Dirac notation and the identity operator, as in blueprint/src/macros/common.tex.
% \providecommand keeps note-local definitions authoritative.
\usepackage{dsfont}
\providecommand{\ket}[1]{|#1\rangle}
\providecommand{\bra}[1]{\langle #1|}
\providecommand{\braket}[2]{\langle #1 | #2 \rangle}
\providecommand{\ketbra}[2]{|#1\rangle\!\langle #2|}
\providecommand{\Id}{\mathds{1}}
\providecommand{\one}{\mathds{1}}
\providecommand{\C}{\mathbb{C}}
\providecommand{\MN}[1]{M_{#1}(\C)}
\providecommand{\MD}{M_D(\C)}

% External referencing for paper sources.  The build helper writes sanitized
% auxiliary files under build/paper-aux/ and excludes duplicated source labels.
\usepackage{xr}

% Import a generated paper auxiliary file with a caller-supplied label prefix.
% The candidate paths support builds from docs/paper-gaps/, the repository root,
% and build/paper-aux/ itself.
\newcommand{\importpaperaux}[2]{%
  \IfFileExists{../../build/paper-aux/#2.aux}{%
    \externaldocument[#1]{../../build/paper-aux/#2}%
  }{%
    \IfFileExists{build/paper-aux/#2.aux}{%
      \externaldocument[#1]{build/paper-aux/#2}%
    }{%
      \IfFileExists{#2.aux}{%
        \externaldocument[#1]{#2}%
      }{}%
    }%
  }%
}

% General external reference macros
\newcommand{\paperref}[2]{\ref{#1-#2}}
\newcommand{\papereqref}[2]{(\ref{#1-#2})}

% CPSV16 source (1606.00608 / MPDO-22-12-17-2.tex) external document and shortcuts
\importpaperaux{cpsv16-}{1606.00608/MPDO-22-12-17-2-tnlean}
\newcommand{\cpsvref}[1]{\paperref{cpsv16}{#1}}
\newcommand{\cpsveqref}[1]{\papereqref{cpsv16}{#1}}

% Verdict marker: \gapnote{<kind>}{<status>}. Kind is the policy
% classification (clarification, local-correction, scope-restriction,
% unfaithful, false-source, open-gap); status is open, wip (elimination
% underway), resolved, or historical. Machine-read by the site index;
% prints a verdict line.
\newcommand{\gapnote}[2]{\par\noindent\textbf{Verdict:} #1 (#2).\par}


\title{Historical Note: Rate-Quantified BNT Cross-Overlap Decay\\
       for a Retired Projection Argument}
\author{}
\date{2026-05-13}

\begin{document}
\maketitle
\gapnote{unfaithful}{historical}

\section{Status}

\textbf{Historical proof archaeology, not an active gap.}
The corrected proportional fundamental theorem is
\leanid{MPSTensor.fundamentalTheorem_proportional_canonicalForm}. It compares
BNT families whose copy weights are nonzero, the standing convention of
CPSV16 line~246 recorded in
\path{docs/paper-gaps/cpsv16_bnt_uniqueness_zero_coefficient.tex}, and obtains a
bijection of the basis sectors with gauge-phase equivalences. No
rate-quantified overlap estimate is needed for the corrected theorem. The
analysis below records the extra hypothesis required by the retired per-block
projection attempt.

\section{Notation}

Fix a sector decomposition $P$ with basis blocks
$\{P.\mathtt{basis}\,j\}_{j=0}^{g-1}$ and spectral levels
$\lambda_j\in\C\setminus\{0\}$ satisfying
$|\lambda_0|>|\lambda_1|>\cdots>|\lambda_{g-1}|$ and
$|\lambda_0|=1$. For physical words $\sigma$ of length $N$, define the
cross-overlap between basis blocks $j$ and $k$ by
\begin{align}
    O_{j,k}(N)
        &:= \mathtt{mpvOverlap}\bigl(P.\mathtt{basis}\,j,
                                    P.\mathtt{basis}\,k\bigr)(N)
    \label{eq:bnt-rate_cross_overlap_definition}\\
        &= \sum_{\sigma}
            \mathtt{mpv}(P.\mathtt{basis}\,j)(\sigma)
            \overline{\mathtt{mpv}(P.\mathtt{basis}\,k)(\sigma)}.
    \label{eq:bnt-rate_cross_overlap_expansion}
\end{align}
This overlap is written \leanid{MPSTensor.mpvOverlap} in the
formalization.\footnote{Declaration \leanid{MPSTensor.mpvOverlap} in
    \doclink{TNLean/MPS/Overlap/Basic}.}

For every ordered pair $j\ne k$, the retired argument sought constants
$C_{j,k}\geq 0$ and rates $\eta_{j,k}\in[0,1)$ such that
\begin{align}
    |O_{j,k}(N)|
        &\leq C_{j,k}\eta_{j,k}^{N}
        \qquad\text{for every }N\in\mathbb N.
    \label{eq:bnt-rate_geometric_cross_overlap_bound}
\end{align}

\section{Cited assertion}

The Step~1 argument in the proof of Theorem~\texttt{thm1} of
Ref.~\cite[line~1182]{Cirac2016MPDO_arXiv} relies on
$O_{j,k_0}(N)\to 0$ for $j\ne k_0$ to discard cross-overlap terms when
projecting an assembled BNT identity onto the distinguished basis block
$V^{(N)}(P.\mathtt{basis}\,k_0)$. The qualitative limit is implicit in the BNT
structure: distinct normal blocks satisfy
\begin{align}
    O_{j,k}(N)
        &\xrightarrow[N\to\infty]{}0
        \qquad\text{for }j\ne k.
    \label{eq:bnt-rate_qualitative_cross_overlap_decay}
\end{align}
The cited argument does not quantify the rate of decay in
(\ref{eq:bnt-rate_qualitative_cross_overlap_decay}).

Definition~4.2 of Ref.~\cite{Cirac2021Matrix} packages
(\ref{eq:bnt-rate_qualitative_cross_overlap_decay}) indirectly through
\emph{eventual linear independence} of the basis MPVs and does not state a
rate. The formalization records this eventual-linear-independence input in the
shared sector-decomposition infrastructure,\footnote{%
    \doclink{TNLean/MPS/SharedInfra/SectorDecomposition}.}
while (\ref{eq:bnt-rate_qualitative_cross_overlap_decay}) is provided
pointwise by
\leanid{MPSTensor.mpvOverlap_tendsto_zero_of_not_gaugePhaseEquiv_cast_left_of_irreducible_TP}.\footnote{%
    \doclink{TNLean/MPS/Overlap/CastDecay}.}
Neither source equips $\eta_{j,k}$ with a quantitative relation to the
spectral-weight ratios.

\section{The retired rate hypothesis}

An earlier formal route introduced an explicit rate hypothesis on a sector
decomposition $P$ and a chosen spectral-level function
$\lambda:\mathrm{Fin}\,P.\mathtt{basisCount}\to\C$. That route has since been
retired together with the per-block-projection layer. The mathematical
hypothesis it tried to isolate required nonnegative families $C$ and $\eta$
such that, for every $j\ne k$,
\begin{align}
    0
        &\leq C_{j,k},
    \label{eq:bnt-rate_nonnegative_prefactor}\\
    0
        &\leq \eta_{j,k}<1,
    \label{eq:bnt-rate_contracting_rate}\\
    \eta_{j,k}\left|\frac{\lambda_j}{\lambda_k}\right|
        &<1,
    \label{eq:bnt-rate_weight_compatibility}
\end{align}
together with the geometric bound
(\ref{eq:bnt-rate_geometric_cross_overlap_bound}) for every $N$.

The surviving \path{TNLean/MPS/FundamentalTheorem/SectorBNT} surface does not package such constants. It
keeps the raw BNT weights in \leanid{MPSTensor.SectorDecomposition} and records
qualitative consequences in
\doclink{TNLean/MPS/FundamentalTheorem/SectorBNT/Api}: the raw power-sum
coefficient identity \leanid{MPSTensor.SectorDecomposition.coeff_eq_sum_weight_pow},
within-family cross-decay
\leanid{MPSTensor.IsBNTCanonicalForm.cross_overlap_basis_tendsto_zero}, the
combined-family linear-independence criterion
\leanid{MPSTensor.IsBNTCanonicalForm.combined_family_eventually_li}, and the
modulus-free non-vanishing coefficient statement
\leanid{MPSTensor.SectorDecomposition.coeff_not_eventually_zero} from
\doclink{TNLean/MPS/SharedInfra/SectorDecomposition}. None of these
declarations supplies the constants in
(\ref{eq:bnt-rate_nonnegative_prefactor})--%
(\ref{eq:bnt-rate_weight_compatibility}).

\section{Where the retired hypothesis was used}

The retired per-block-projection route carried an abstract hypothesis excluding
asymptotic cancellation of the assembled overlap against a chosen projection
block $k_0$. The obstruction analysis in
\path{docs/audits/2026-05-13_cpsv16_ft_path_beta_pr2_5_design.md} shows that,
for a non-dominant $k_0$, the argument first isolates
$V^{(N)}(P.\mathtt{basis}\,k_0)$ on the right and then renormalizes by
$\lambda_{k_0}^N$. Every $j\ne k_0$ must consequently contribute a term
satisfying
\begin{align}
    |O_{j,k_0}(N)|
    \left|\frac{\lambda_j}{\lambda_{k_0}}\right|^N
        &\leq
        C_{j,k_0}\eta_{j,k_0}^N
        \left|\frac{\lambda_j}{\lambda_{k_0}}\right|^N
    \label{eq:bnt-rate_weighted_overlap_first_bound}\\
        &=
        C_{j,k_0}
        \left(\eta_{j,k_0}
            \left|\frac{\lambda_j}{\lambda_{k_0}}\right|\right)^N.
    \label{eq:bnt-rate_weighted_overlap_geometric_bound}
\end{align}
The right-hand side of
(\ref{eq:bnt-rate_weighted_overlap_geometric_bound}) converges to zero exactly
when (\ref{eq:bnt-rate_weight_compatibility}) holds. The qualitative limit
(\ref{eq:bnt-rate_qualitative_cross_overlap_decay}) alone provides no
constants $C_{j,k_0}$ and $\eta_{j,k_0}<1$, and hence no geometric rate. Such a
bound requires a separate finite-dimensional spectral-radius analysis of the
mixed transfer operator. Even after obtaining a geometric rate, the comparison
in (\ref{eq:bnt-rate_weight_compatibility}) remains an additional compatibility
condition required by the retired argument.

\section{Classification}

\emph{Historical analytic obstruction.} The rate condition
(\ref{eq:bnt-rate_weight_compatibility}), together with
(\ref{eq:bnt-rate_geometric_cross_overlap_bound}), is a structural input beyond
what Ref.~\cite[\S II]{Cirac2016MPDO_arXiv} and
Definition~4.2 of Ref.~\cite{Cirac2021Matrix} make explicit. Obtaining a
geometric rate requires a separate finite-dimensional spectral-radius input
for the mixed transfer operator underlying
\leanid{MPSTensor.mpvOverlap_tendsto_zero_of_not_gaugePhaseEquiv_cast_left_of_irreducible_TP}.
Moreover, the compatibility inequality
$\eta_{j,k}|\lambda_j/\lambda_k|<1$ is not implied by the BNT structure. The
retired projection argument would therefore have needed this condition as an
additional hypothesis or a separate analytic theorem.

The design memo
\path{docs/audits/2026-05-13_cpsv16_ft_path_beta_pr2_5_design.md} records the
abandoned spectral-gap proposal. The corrected \path{TNLean/MPS/FundamentalTheorem/SectorBNT} theorem uses
coefficient comparison and combined-family linear independence instead. This
note documents why the rate-quantified projection route must not be revived as
a purported obligation of the proportional fundamental theorem.

\bibliographystyle{alpha}
\bibliography{references}

\end{document}
